\documentclass[10pt, aspectratio=169]{beamer}
\usetheme{Madrid}
\usecolortheme{whale}

% Força o uso de fontes serifadas padrão para matemática e evita bugs no MiKTeX
\usefonttheme[onlymath]{serif}

\usepackage[utf8]{inputenc}
\usepackage[brazil]{babel}
\usepackage{booktabs}
\usepackage{amsmath}
\usepackage{amsfonts}

% --- Configurações Iniciais ---
\title[RNAs vs Sistemas Fuzzy]{Análise Comparativa de Desempenho: Redes Neurais Artificiais vs. Sistemas Fuzzy}
\subtitle{Cenários Complexos de Classificação e Regressão (UCI Repository)}
\author[G. Morcatti \and J. Lima]{Geovane de Freitas Queiroz Morcatti \and João Paulo de Lima}
\institute[CEFET-MG]{Centro Federal de Educação Tecnológica de Minas Gerais (CEFET-MG) \\ Mestrado em Modelagem Matemática e Computacional}
\date[Amanhã]{Julho de 2026}

\begin{document}

% --- SLIDE 1: CAPA ---
\begin{frame}
    \titlepage
\end{frame}

% --- SLIDE 2: METODOLOGIA ---
\begin{frame}{Metodologia e Rigor Experimental}
    \begin{itemize}
        \item \textbf{Divisão Isonômica dos Dados:} Proporção estrita de 60\% Treinamento, 20\% Validação e 20\% Teste[cite: 74].
        \item \textbf{Mitigação da Estocasticidade:} Protocolo rigoroso de \textbf{21 execuções independentes} por algoritmo, avaliando média e desvio-padrão ($\pm \sigma$)[cite: 75].
        \item \textbf{Engenharia de Recursos:} Uso de \textit{SelectKBest} (ANOVA) para reduzir o espaço a 2 atributos nos modelos Fuzzy, viabilizando o universo de regras lógicas[cite: 76].
    \end{itemize}
\end{frame}

% --- SLIDE 3: DATASETS ---
\begin{frame}{Conjuntos de Dados Avaliados (UCI Repository)}
    \centering
    \small
    \begin{tabular}{lcccc}
        \toprule
        \textbf{Dataset} & \textbf{Tarefa} & \textbf{Amostras} & \textbf{Atributos} & \textbf{Target} \\
        \midrule
        1. Iris & Classif. Multi. & 150 & 4 Num. & 3 Classes (Balanceado) [cite: 77, 78] \\
        2. Breast Cancer & Classif. Bin. & 569 & 30 Num. & M/B (Alta Dimensionalidade) [cite: 79] \\
        3. Heart Disease & Classif. Bin. & 270 & 13 Mistos & S/N (Dados Híbridos) [cite: 80] \\
        4. Wine Quality & Regressão & 1.599 & 11 Num. & Contínuo (Notas 3 a 8) [cite: 81] \\
        \bottomrule
    \end{tabular}
\end{frame}

% --- SLIDE 4: RESULTADOS CLASSIFICAÇÃO ---
\begin{frame}{Resultados — Tarefas de Classificação}
    \centering
    \textbf{Acurácia Média $\pm$ Desvio-Padrão ($\sigma$)} \\
    \vspace{0.15cm}
    \scriptsize
    \begin{tabular}{lccccc}
        \toprule
        \textbf{Dataset} & \textbf{MLP Rasa (A)} & \textbf{MLP Prof. (B)} & \textbf{Rede Neural RBF} & \textbf{Mamdani A (Tri.)} & \textbf{Mamdani B (Gau.)} \\
        \midrule
        Iris & $0.9571 \pm 0.0426$ & $\mathbf{0.9667 \pm 0.0341}$ & $0.9333 \pm 0.0436$ & $0.6206 \pm 0.0281$ & $0.3349 \pm 0.0071$ \\
        Breast Cancer & $\mathbf{0.9758 \pm 0.0130}$ & $0.9683 \pm 0.0187$ & $0.9449 \pm 0.0180$ & $0.1157 \pm 0.0252$ & $0.1257 \pm 0.0215$ \\
        Heart Disease & $0.8104 \pm 0.0450$ & $0.8051 \pm 0.0445$ & $\mathbf{0.8201 \pm 0.0388}$ & $0.7381 \pm 0.0376$ & -- \\
        \bottomrule
    \end{tabular}
    \vspace{0.15cm}
    \begin{block}{Destaque do Experimento}
        A \textbf{Rede RBF superou as MLPs na base Heart Disease (82.01\%)}[cite: 83]. O modelo Mamdani sofreu forte impacto de \textit{underfitting} no \textit{Breast Cancer} devido à redução drástica para 2 atributos (*Maldição da Dimensionalidade*)[cite: 89].
    \end{block}
\end{frame}

% --- SLIDE 5: RESULTADOS REGRESSÃO ---
\begin{frame}{Resultados — Tarefa de Regressão (Wine Quality)}
    \centering
    \small
    \begin{tabular}{lccc}
        \toprule
        \textbf{Algoritmo Avaliado} & \textbf{MSE Médio} & \textbf{MAE Médio} & \textbf{$R^2$ Médio} \\
        \midrule
        MLP Rasa (Config. A) & $0.5533$ & $0.5693$ & $0.1757$ [cite: 85] \\
        MLP Profunda (Config. B) & $\mathbf{0.4415}$ & $\mathbf{0.5124}$ & $\mathbf{0.3410}$ [cite: 85] \\
        Rede Neural RBF & $10.2308$ & $3.1076$ & $-14.2824$ [cite: 86] \\
        Fuzzy Takagi-Sugeno (T-S) & $2.2396$ & $1.2573$ & $-2.3428$ [cite: 86] \\
        \bottomrule
    \end{tabular}
    \vspace{0.2cm}
    \begin{alertblock}{Análise de Fronteiras e Erros}
        O $R^2$ negativo da RBF decorre do modelo \textit{RbfClassifier} adaptado[cite: 86]. Tentar prever um espaço contínuo mapeando fronteiras discretas gera aproximações em degraus, elevando o erro quadrático residual[cite: 87].
    \end{alertblock}
\end{frame}

% --- SLIDE 6: CONCLUSÃO ---
\begin{frame}{Discussão Crítica e Conclusões}
    \begin{itemize}
        \item \textbf{Trade-off da Inteligência Computacional:} Redes Neurais (MLP e RBF) garantem superioridade preditiva matemática global, agindo como caixas-pretas[cite: 88].
        \item \textbf{Estabilidade Fuzzy:} Embora penalizados na acurácia pela perda de atributos essenciais, os modelos Fuzzy entregam os \textbf{menores desvios-padrão} do teste e permitem auditoria semântica clara[cite: 89].
        \item \textbf{Conclusão:} A escolha do modelo ideal deve balancear a necessidade de precisão pura com o nível de explicabilidade exigido pelo domínio prático[cite: 90].
    \end{itemize}
\end{frame}

\end{document}